%%%%%%%%%%%%%%%%%%%%%%%%%%%%%%%%%%%%%%%%
\chapter{Objetivos}
%%%%%%%%%%%%%%%%%%%%%%%%%%%%%%%%%%%%%%%%

El presente capítulo define el objetivo general y los objetivos específicos que orientan esta revisión sistemática, estableciendo el alcance y los productos esperados del trabajo.

%%%%%%%%%%%%%%%%%%%%%%%%%%%%%%%%%%%%%%%%
\section{Objetivo general}
\begin{quote}
Comparar las técnicas de \textit{fine-tuning} para \textit{Large Language Models} mediante revisión sistemática PRISMA, identificando relaciones entre características técnicas, rendimiento en distintos dominios de aplicación, y patrones de aplicabilidad reportados en la literatura académica entre 2020 y 2025.
\end{quote}

%%%%%%%%%%%%%%%%%%%%%%%%%%%%%%%%%%%%%%%%
\section{Objetivos específicos}

Para alcanzar el objetivo general, se plantean tres objetivos específicos que cubren las etapas de sistematización, análisis y síntesis de la evidencia empírica.

\begin{itemize}
    \item Sistematizar la literatura sobre \textit{fine-tuning} de \textit{Large Language Models} publicada entre 2020 y 2025 aplicando metodología PRISMA, extrayendo datos sobre características técnicas, configuraciones experimentales, métricas de evaluación y resultados de rendimiento en distintos dominios de aplicación.

    \item Analizar el rendimiento de las técnicas de \textit{fine-tuning} según dominio de aplicación, estableciendo \textit{trade-offs} y patrones de aplicabilidad reportados en la literatura.

    \item Sintetizar criterios de selección de técnicas según características del contexto operativo (disponibilidad de datos, dominio de aplicación, objetivos de rendimiento), identificando limitaciones reportadas y direcciones futuras de investigación.
\end{itemize}